% Generated draft: requires mathematical and source review.
\documentclass[11pt]{article}
\usepackage[margin=25mm]{geometry}
\usepackage{fontspec}
\setmainfont{Latin Modern Roman}
\usepackage{amsmath,amssymb,mathtools,unicode-math}
\setmathfont{Latin Modern Math}
\usepackage{enumitem,needspace,xurl,hyperref}
\setlength{\parindent}{0pt}
\setlength{\parskip}{5pt}
\setlength{\emergencystretch}{3em}
\newcommand{\N}{\mathbb{N}}
\newcommand{\Z}{\mathbb{Z}}
\newcommand{\Q}{\mathbb{Q}}
\newcommand{\R}{\mathbb{R}}
\newcommand{\C}{\mathbb{C}}
\newcommand{\F}{\mathbb{F}}
\newcommand{\A}{\mathbb{A}}
\newcommand{\PP}{\mathbb{P}}
\newcommand{\E}{\mathbb{E}}
\newcommand{\eps}{\varepsilon}
\newcommand{\dd}{\,\mathrm{d}}
\begin{document}
\section*{UnsolvedMath ID: 800423}
\textbf{Source problem code:} catalogue entry
\subsection{Definitions and mathematical statement}
Let \(K\) be a number field, let \(\overline K\) be an algebraic closure, and write \(G_K=\operatorname{Gal}(\overline K/K)\). Let \(f:\PP^N\to\PP^N\) be an endomorphism defined over \(K\), of algebraic degree \(d\ge 2\), meaning \(f^*\mathcal O_{\PP^N}(1)\cong\mathcal O_{\PP^N}(d)\). Write \(f^{\circ n}\) for its \(n\)-fold iterate. A point \(P\in\PP^N(K)\) is periodic if \(f^{\circ m}(P)=P\) for some integer \(m\ge 1\).

For each integer \(n\ge 1\), put
\[
X_n(f,P)=\{Q\in\PP^N(\overline K):f^{\circ n}(Q)=P\},
\qquad
r_n(f,P;K)=\#\bigl(G_K\backslash X_n(f,P)\bigr).
\]
The set \(X_n(f,P)\) is finite and is preserved by \(G_K\). In this formulation, \((f,P)\) is \emph{eventually stable over \(K\)} if
\[
\exists C\in\Z_{\ge 1}\ \forall n\in\Z_{\ge 1},\qquad r_n(f,P;K)\le C.
\]
This counts Galois orbits of distinct geometric inverse images, without ramification multiplicities. The supplied excerpt does not define eventual stability, so this convention requires confirmation against the cited source. It also leaves the range of the projective dimension \(N\) implicit.

\textbf{Conjecture 46 (supplied formulation).} For every number field \(K\), every endomorphism \(f:\PP^N\to\PP^N\) defined over \(K\) of degree at least \(2\), and every \(P\in\PP^N(K)\) satisfying
\[
\forall m\in\Z_{\ge 1},\qquad f^{\circ m}(P)\ne P,
\]
the pair \((f,P)\) is eventually stable over \(K\). The bound \(C\) may depend on \(K,f,P\); no uniform bound across these data is asserted. The hypothesis excludes periodic points but permits strictly preperiodic points.
\subsection{Short English statement}
The conjecture says that, for a nonperiodic rational point of a projective endomorphism over a number field, the number of Galois orbits among its depth-\(n\) inverse images remains bounded as \(n\) grows. The supplied excerpt leaves the definition of eventual stability implicit; the orbit-counting convention used here needs source confirmation.
\subsection{Sources}
\begin{itemize}
\item Supplied local UnsolvedMath registry record, permanent ID \(800423\), exact source code \texttt{catalogue entry}, titled \emph{Open Problems in Arithmetic Dynamics --- Conjecture 46}. The registry attributes the problem to Benedetto, DeMarco, Ingram, Jones et al., with proposed year \(2018\).
\item Benedetto, DeMarco, Ingram, Jones et al., \emph{Current Trends and Open Problems in Arithmetic Dynamics} (2018), Conjecture 46. Supplied source URL: \url{https://arxiv.org/abs/1806.04980}. The record reports access on 2026-10-16 and extraction by \texttt{arxiv\_latex}; these are supplied provenance statements, not independent source verification.
\item The registry labels the problem \texttt{open} and records an update dated 2026-10-16. This is supplied dated curation, not an independently verified assessment of mathematical or literature status. Supplied upstream revision: \texttt{d371ee917ac9647fbb06cd6fc529deec6f3dda82}.
\end{itemize}
\end{document}
